% Section 18.10, from lectures/L18/README.md.

\section{Summary}
\label{c18:sec:review}

\needspace{6\baselineskip}
\subsection*{What you should be able to do now}
\begin{itemize}
  \item Explain the difference between a pulse and a sticky level, and why the register layer has to
    convert between them.
  \item Implement \file{register_bank.vhd} per the register map: the STATUS bits' setting and
    clearing, \code{TX_SEND} as a one-cycle pulse, and the masking on write.
  \item Read an SPI timing diagram in mode 0 and say which bit is on MOSI at every edge.
  \item Explain why an SPI slave in an FPGA oversamples SCK instead of using it as a clock.
  \item Read and use \file{spi_slave.vhd} without writing it.
\end{itemize}

\needspace{6\baselineskip}
\subsection*{Questions to test yourself with}
\begin{enumerate}
  \item Why can a driver polling every five microseconds not see a 20~ns pulse?
  \item What sets STATUS bit 0, and what clears it?
  \item Why does the error bit latch on a \emph{rising edge} of \code{error} and not on the level?
  \item A second frame is received before the first has been acknowledged. What happens, and why is
    that a deliberate choice rather than a bug?
  \item STATUS bit 0 is set by three things during operation, besides reset. Which, and why is
    \code{tx_done} alone not enough?
  \item What does \code{TX_ABORT} do, and what does it \emph{not} reach?
  \item A write of \code{0xFFFFFFFF} to \code{TX_DLC} reads back as \code{0x0000000F}. Where does the
    masking happen, and why does the bank not validate that the value is at most 8?
  \item Why does \code{spi_slave} not use SCK as a clock? What would go wrong if it did?
  \item Why does MISO change value on SCK's falling edge?
\end{enumerate}

\needspace{6\baselineskip}
\subsection*{Target}
\code{register_bank} is a self-contained module with a complete testbench and should be finishable
during the session and the week after. It does not depend on SPI at all, so it can be written in
parallel with somebody else in the group reading up on the transaction format. Put the
\file{spi_slave.vhd} handed out into the repository unchanged, and read through
\file{spi_reg_bridge_tb.vhd} in preparation for the next chapter: it is the contract for the module
you write then.

\needspace{6\baselineskip}
\subsection*{Looking ahead}
\Chapref{c19:ch} builds the SPI bridge, the top level \code{can_spi_node}, and takes the whole chain
out onto real hardware. It is the project's last session.
